\documentclass[11pt,a4paper]{coverletter}

\senderblock
  {Radheem Bin Razi}
  {Distributed Systems \& Cloud-Native Engineer}
  {Ilmenau, Germany \textbar\ \href{mailto:sheikh.radheem@gmail.com}{sheikh.radheem@gmail.com} \textbar\ \href{https://radheem.github.io/my-cv}{portfolio}}

\begin{document}

%% ===== ENGLISH =====
\recipient{BMW Group}{Hiring Team}{}

\opening{Dear Hiring Team,}

Working at the seam between software and large-scale data is where I do my best work, so a working-student role at BMW Group — where software increasingly defines the car and the business runs serious data, cloud, and AI platforms — was one I genuinely wanted to apply for. I am a Master of Research student at TU Ilmenau with five years of professional software-engineering experience, and I would bring hands-on backend, cloud-native, and data/ML skills to a BMW team a few days a week.

Most of my work has been building and operating distributed systems. On a stealth platform I built Go microservices over gRPC and NATS JetStream with a best-effort ETL pipeline into PostgreSQL and Dgraph, all deployed on Kubernetes — the kind of reliable, observable backend that production scale demands. Earlier, as a design authority at a high-throughput exchange, I reviewed code and architecture and drove improvements where the system was weakest, working day to day with Kafka, PostgreSQL, Docker, and Terraform.

Data and ML platforms are a particular interest. For my research I built an end-to-end AI/ML framework on Kubernetes (Kubeflow training pipelines, KServe serving) driven by a config-driven Python SDK, wired to InfluxDB and a feature store and graded 1.0 — exactly the kind of platform-shaped work a large software organisation invests in. I default to clean, tested code and treat observability and reviews as how a team gets better, and I lean on modern tooling like RAG and MCP in my own projects; BMW's scale is where I would like to apply and deepen that.

As an enrolled M.Sc. student in Germany I am eligible to work as a working student — around 20 hours per week during the lecture period and full-time in the breaks. I am based in Ilmenau, already authorized to work here, and happy with a hybrid arrangement or relocating within Germany. My German is at A2 and improving; my working language is English. Thank you for your time and consideration — I would welcome the chance to contribute to a BMW software or data team.

\closing{Sincerely,}{Radheem Bin Razi}

\clearpage
\selectlanguage{ngerman}

%% ===== DEUTSCH =====
\recipient{BMW Group}{Personalabteilung}{}

\opening{Sehr geehrtes Hiring-Team,}

An der Schnittstelle zwischen Software und großvolumigen Daten leiste ich meine beste Arbeit, weshalb eine Werkstudentenstelle bei der BMW Group — wo Software zunehmend das Fahrzeug prägt und das Unternehmen anspruchsvolle Daten-, Cloud- und KI-Plattformen betreibt — eine Stelle war, auf die ich mich wirklich bewerben wollte. Ich bin Master-of-Research-Student an der TU Ilmenau mit fünf Jahren beruflicher Erfahrung in der Softwareentwicklung und würde an einigen Tagen pro Woche praxiserprobte Kenntnisse in Backend, Cloud-Native sowie Daten/ML in ein BMW-Team einbringen.

Ein Großteil meiner Arbeit bestand im Aufbau und Betrieb verteilter Systeme. Auf einer Stealth-Plattform habe ich Go-Microservices über gRPC und NATS JetStream mit einer Best-Effort-ETL-Pipeline nach PostgreSQL und Dgraph entwickelt, alles auf Kubernetes bereitgestellt — genau jene zuverlässige, beobachtbare Backend-Architektur, die produktive Skalierung erfordert. Zuvor habe ich als Design-Verantwortlicher bei einer Börse mit hohem Durchsatz Code und Architektur überprüft und Verbesserungen dort vorangetrieben, wo das System am schwächsten war; täglich arbeitete ich dabei mit Kafka, PostgreSQL, Docker und Terraform.

Daten- und ML-Plattformen sind ein besonderes Interessensgebiet von mir. Im Rahmen meiner Forschung habe ich ein durchgängiges KI-/ML-Framework auf Kubernetes (Kubeflow-Trainings-Pipelines, KServe-Serving) entwickelt, gesteuert über ein konfigurationsgetriebenes Python-SDK, angebunden an InfluxDB und einen Feature Store und mit der Note 1,0 bewertet — genau die plattformartige Arbeit, in die eine große Softwareorganisation investiert. Ich setze standardmäßig auf sauberen, getesteten Code und betrachte Observability und Reviews als die Art und Weise, wie ein Team besser wird; in meinen eigenen Projekten nutze ich moderne Werkzeuge wie RAG und MCP. Bei der Größenordnung von BMW würde ich diese Fähigkeiten gerne anwenden und vertiefen.

Als immatrikulierter M.Sc.-Student in Deutschland bin ich berechtigt, als Werkstudent zu arbeiten — rund 20 Stunden pro Woche während der Vorlesungszeit und in vollem Umfang in den Semesterferien. Ich bin in Ilmenau ansässig, bereits zur Arbeit hier berechtigt und gerne zu einem hybriden Arbeitsmodell oder einem Umzug innerhalb Deutschlands bereit. Mein Deutsch liegt auf A2-Niveau und verbessert sich stetig; meine Arbeitssprache ist Englisch. Vielen Dank für Ihre Zeit und Ihre Berücksichtigung — ich würde mich freuen, zu einem Software- oder Datenteam von BMW beitragen zu dürfen.

\closing{Mit freundlichen Grüßen,}{Radheem Bin Razi}

\end{document}
